\documentclass[11pt,a4paper]{article}
\usepackage{amsmath,amssymb,geometry,graphicx,booktabs,hyperref,microtype}
\geometry{margin=2.5cm}
\hypersetup{colorlinks=true,linkcolor=blue,citecolor=blue,urlcolor=blue}

\title{\textbf{Paper 47: C$_2$ Coordinate Validation}\\
\large $\delta$-Sweep and Gate--Adversarial Interaction}
\author{VCSM Project}
\date{2026-03-10}

\begin{document}
\maketitle

\begin{abstract}
We test the two remaining untested coordinates of the C$_2$ causal
reparametrisation introduced in Paper 46.  Experiment~A sweeps the causal depth
ratio $\delta = W_\text{zone}/r_\text{wave}$ across five zone configurations
while holding $\Omega = \mathrm{WR}\cdot\phi_w$ fixed at 62.4.  Zone encoding
quality, measured by the non-adjacent/adjacent distance ratio, peaks at
$\delta = 5$ (standard 4-zone grid) and degrades toward unity as
$\delta \to 2$, confirming that $\delta \geq r_\text{wave}$ is required for
coherent zone-level gradients.  Experiment~B crosses gate condition
($\mathrm{SS} \in \{0,5,10,15,20,\text{rand}_{p=0.60}\}$) with adversarial
wave amplitude ($a \in \{0, 0.125, 0.25, 0.50\}$) during a post-encoding
adversarial phase.  An unexpected finding emerges: at encoding checkpoint
$T = 2000$ (near the commitment epoch $t^* \approx 2000$ from Paper~26), the
cosine fidelity of the fieldM snapshot is near zero or negative even under
normal waves, revealing that the metric captures \emph{commitment rate} rather
than adversarial resistance at pre-commitment snapshots.  High-SS gates slow
commitment; low-SS and random gates accelerate it.  Together the two experiments
tighten the C$_2$ picture: $\delta$ is the third operative coordinate, and gate
selectivity interacts with the commitment epoch, not only with adversarial noise.
\end{abstract}

\section{Background}

Paper~46 established the C$_2$ causal reparametrisation of the VCSM parameter
space.  Two of its three coordinates were experimentally validated:
\begin{itemize}
  \item $\Omega = \mathrm{WR}\cdot\phi_w(r_\text{wave})$ --- causal dose.
        Paper~46 Exp~B confirmed: fixing $\Omega$ makes $\mathrm{sg4}$ invariant
        to $r_\text{wave}$ (CV $= 0.127$).
  \item $\beta_s$ --- seeding scale.  Held fixed throughout the paper series.
\end{itemize}
The third coordinate, $\delta = W_\text{zone}/r_\text{wave}$, had not been
directly swept.  Intuitively, $\delta$ controls whether wave events are
\emph{zone-local} ($\delta \gg 1$) or \emph{cross-zone} ($\delta \approx 1$).
Paper~47 Experiment~A fills this gap.

Paper~26 showed that the zone-mean field vector $\mathbf{F}(t)$ is
anti-correlated with its final value for $t < t^* \approx 1600$--$2000$: the
system encodes the \emph{wrong} zone pattern before committing.
Paper~47 Experiment~B tests gate--adversarial interactions; a snapshot at
$T = 2000$ near $t^*$ creates an interpretive subtlety that itself becomes a
finding.

\section{Experiments}

\subsection{Experiment A: $\delta$-sweep}

\textbf{Design.}  Five zone configurations:
$N_\text{zones} \in \{2,4,5,8,10\}$,
giving $\delta \in \{10, 5, 4, 2.5, 2\}$ with $r_\text{wave}=2$ fixed.
All other C$_2$ coordinates held constant: $\Omega = 62.4$ (WR$=4.8$,
$r_\text{wave}=2$), $\beta_s = 0.25$, $\mathrm{SS}=10$.  Five seeds,
$T = 3000$, tail average over last 30 samples.

\textbf{Primary metric.} Raw $\mathrm{sg4}$ (mean pairwise L2 between zone-mean
fieldM vectors) and the non-adjacent/adjacent ratio
$R_\text{na} = \langle d_\text{nonadj}\rangle / \langle d_\text{adj}\rangle$.
$R_\text{na} > 1$ indicates monotonic position encoding (zones that are farther
apart in space have more different fieldM).  $R_\text{na} < 1$ indicates that
adjacent zones are \emph{more} different than non-adjacent zones, a sign of
sub-zone or alternating structure.

\textbf{Results.}

\begin{table}[h]
\centering
\caption{Experiment A: $\delta$-sweep results (mean $\pm$ SE over 5 seeds).}
\label{tab:expa}
\begin{tabular}{rrrrrr}
\toprule
$N_\text{zones}$ & $\delta$ & $W_\text{zone}$ & sg4 & sg4 SE & $R_\text{na}$ \\
\midrule
 2 & 10.0 & 20 & 0.0174 & 0.0013 & n/a \\
 4 &  5.0 & 10 & 0.0244 & 0.0020 & 1.061 \\
 5 &  4.0 &  8 & 0.0217 & 0.0014 & 0.971 \\
 8 &  2.5 &  5 & 0.0289 & 0.0010 & 0.971 \\
10 &  2.0 &  4 & 0.0327 & 0.0015 & 1.036 \\
\bottomrule
\end{tabular}
\end{table}

Raw sg4 increases monotonically as $\delta$ decreases.  This is partly a
counting artifact: with $N_\text{zones}=10$ there are $\binom{10}{2}=45$
zone pairs versus $\binom{4}{2}=6$, and each pair involves only 4 cells of
fieldM, increasing per-pair variance.  The encoding-quality signal is
$R_\text{na}$:

\begin{itemize}
  \item $\delta = 5$ ($N_\text{zones}=4$, standard configuration):
        $R_\text{na} = 1.061$.  Clear monotonic position gradient.
  \item $\delta = 4$ ($N_\text{zones}=5$):
        $R_\text{na} = 0.971 < 1$.  The zone structure has degraded;
        some adjacent zone pairs are already more distinct than non-adjacent.
  \item $\delta = 2.5$ ($N_\text{zones}=8$):
        $R_\text{na} = 0.971$.  Consistent with $\delta=4$; wave events
        routinely span zone boundaries.
  \item $\delta = 2.0$ ($N_\text{zones}=10$):
        $R_\text{na} = 1.036$.  Slight recovery; at this scale the even/odd
        wave polarity pattern (suppression vs excitation) reinstates a
        two-class structure at the wave-level grain.
  \item $\delta = 10$ ($N_\text{zones}=2$):
        $R_\text{na}$ undefined (only one pair, which is adjacent).
\end{itemize}

\textbf{Interpretation.}  $\delta$ is an operative coordinate.  The viable
range for zone-level encoding is $\delta \geq 5$.  At $\delta < 5$, the wave
footprint ($r_\text{wave}=2$, diameter~5) exceeds or matches the zone width
($W_\text{zone}=4$--$8$), and the copy-forward loop builds sub-zone spatial
gradients rather than zone-level ones.  The raw sg4 increase at small $\delta$
reflects finer-grained spatial differentiation, but not the
\emph{hierarchical} zone structure needed for reliable zone-identity encoding.

The $R_\text{na}$ dip at $\delta=4$ versus recovery at $\delta=2$ is consistent
with an interference pattern: at $\delta=4$ zone boundaries fall mid-wave,
creating ambiguous overlap; at $\delta=2$ zones are narrow enough that the
even/odd polarity rule consistently assigns one suppression class and one
excitation class per wave event, re-imposing a two-class distinction at the
sub-zone level.

\subsection{Experiment B: Gate $\times$ Adversarial Pressure}

\textbf{Design.}  Two-phase protocol:
\begin{enumerate}
  \item \textbf{Encoding phase} [$T = 0$\ldots$2000$]: Full VCML with SS$=10$,
        standard waves.  All conditions identical.  Builds zone structure.
  \item \textbf{Test phase} [$T = 2000$\ldots$3000$]: Gate switches to condition-
        specific SS or $\text{rand}_{p=0.60}$.  Adversarial waves (polarity
        flipped, amplitude scaled by $a \in \{0, 0.125, 0.25, 0.50\}$).
\end{enumerate}
Primary metric: \emph{fidelity} = mean cosine similarity between zone-mean
fieldM at the end of the test phase and the baseline snapshot taken at $T=2000$.
$4\ \text{adv\_amp} \times 6\ \text{gates} \times 5\ \text{seeds} = 120$ runs.

\textbf{Results.}

\begin{table}[h]
\centering
\caption{Experiment B: final fidelity (mean over 5 seeds) at each
         adv\_amp $\times$ gate combination.}
\label{tab:expb}
\begin{tabular}{lrrrrrr}
\toprule
adv\_amp & SS=0 & SS=5 & SS=10 & SS=15 & SS=20 & rand$_{0.60}$ \\
\midrule
0.000 & $-0.219$ & $ 0.052$ & $ 0.107$ & $ 0.172$ & $ 0.089$ & $-0.144$ \\
0.125 & $-0.144$ & $-0.066$ & $-0.068$ & $-0.061$ & $-0.076$ & $-0.053$ \\
0.250 & $-0.070$ & $-0.081$ & $-0.108$ & $-0.112$ & $-0.058$ & $ 0.021$ \\
0.500 & $-0.212$ & $-0.077$ & $-0.157$ & $-0.213$ & $-0.281$ & $-0.102$ \\
\bottomrule
\end{tabular}
\end{table}

A striking result: even with \emph{no adversarial input} ($a=0$), nearly all
conditions show fidelity $\leq 0.17$.  SS$=0$ reaches $-0.22$: the fieldM has
inverted polarity relative to the $T=2000$ snapshot.  This is not a failure of
the system but a reflection of the commitment epoch.

\textbf{Commitment epoch interpretation.}  Paper~26 established that the
zone-mean cosine similarity $C(t_\text{ref}, T_\text{end})$ is \emph{negative}
for $t_\text{ref} < t^* \approx 1600$--$2000$: the fieldM encodes an
anti-correlated pre-pattern before committing to its final state.  A baseline
snapshot at $T=2000$ may be at or just before commitment for some seeds.

Under this interpretation:
\begin{itemize}
  \item \textbf{Low-SS gates (SS$=0$, rand$_{0.60}$)}: high write rate
        accelerates commitment; fieldM rapidly crosses the sign flip and ends
        at the true attractor, which is anti-correlated with the $T=2000$
        pre-committed snapshot.  Fidelity becomes strongly negative.
  \item \textbf{High-SS gates (SS$=15$)}: low write rate slows commitment;
        fieldM remains near the $T=2000$ state.  Fidelity stays positive.
  \item \textbf{Adversarial input} ($a > 0$): disrupts commitment by driving
        mid\_mem in the wrong direction, preventing any gate from accumulating
        a clean committed state.  All fidelity values converge toward zero or
        slightly negative.
\end{itemize}

The ordering at $a=0$ --- SS$=15 >$ SS$=10 >$ SS$=5 >$ SS$=20 >$
rand$_{0.60} >$ SS$=0$ --- maps onto commitment rate: SS$=15$ is near the
optimal gate for resisting the sign flip of the pre-commitment trajectory.
SS$=20$ (tightest gate) is slightly lower, possibly because too few writes
prevent even the minimum consolidation needed to maintain the $T=2000$ state.

At $a=0.25$, rand$_{0.60}$ achieves fidelity $+0.021$ (best), consistent with
the Paper~46 adversarial filter result: random writes occasionally catch the
high-signal onset of adversarial waves but do not systematically consolidate
adversarial mid\_mem.

\textbf{Experimental limitation.}  The $T=2000$ snapshot is within the
commitment epoch window.  A cleaner test of adversarial resistance would use
$T_\text{encode} = 3000$ (as in Paper~44), where fieldM is past $t^*$ and the
baseline is the fully committed state.  The present experiment inadvertently
probes the interaction between gate selectivity and commitment dynamics, which
is itself a non-trivial finding but distinct from the intended design.

\section{Discussion}

\subsection{$\delta$ as Third C$_2$ Coordinate}

Table~\ref{tab:expa} confirms $\delta$ is operative: encoding quality
($R_\text{na}$) peaks at $\delta = 5$ and degrades at smaller values.  The
condition $\delta \geq r_\text{wave} = 2$ is necessary but not sufficient;
$\delta \geq 5$ appears to be the practical minimum for zone-level gradients
with $r_\text{wave}=2$.

This is consistent with the copy-forward loop mechanism.  A cell at the zone
boundary receives wave events from both its own zone and the neighboring zone.
For zone-level encoding to dominate, the majority of a cell's wave events must
originate from one zone.  The fraction of events from its home zone is
approximately $\delta/(\delta+1)$: at $\delta=5$, 83\% home-zone events; at
$\delta=2$, 67\%; at $\delta=1$, 50\% (random assignment).  The observed
$R_\text{na}$ degradation between $\delta=5$ and $\delta=4$ is consistent with
this crossing the practical discrimination threshold.

\subsection{Gate Selectivity and the Commitment Epoch}

The Experiment~B result unifies three prior findings:
\begin{enumerate}
  \item \textbf{Paper~26}: zone structure anti-correlates with its final state
        for $t < t^*$.
  \item \textbf{Paper~46}: the quiescence gate is an adversarial filter (high
        SS protects against encoding mid-perturbation signals).
  \item \textbf{Paper~47 Exp~B}: at $T_\text{encode} = t^*$, gate selectivity
        controls commitment rate: conservative gates slow the sign flip, liberal
        gates accelerate it.
\end{enumerate}
These are three aspects of a single mechanism: the gate controls \emph{how fast
fieldM tracks mid\_mem}, and mid\_mem itself passes through a sign flip
during the commitment epoch.  A high-SS gate applies a temporal low-pass filter
that smooths over this flip; a low-SS gate allows fieldM to follow mid\_mem
through the transition instantly.

This has implications for temporal encoding: the optimal SS is not a fixed
parameter but a function of the adversarial pressure level \emph{and} the
distance from the commitment epoch.

\subsection{C$_2$ Coordinates: Summary Status}

\begin{table}[h]
\centering
\caption{C$_2$ causal reparametrisation: validation status.}
\label{tab:c2}
\begin{tabular}{lllll}
\toprule
Coordinate & Definition & Prediction & Paper & Status \\
\midrule
$\Omega$ & $\mathrm{WR}\cdot\phi_w(r)$ & flat sg4 at fixed $\Omega$ &
  P46 Exp B & CONFIRMED (CV=0.127) \\
$\delta$ & $W_\text{zone}/r_\text{wave}$ & sg4 quality peaks at $\delta \approx 5$ &
  P47 Exp A & CONFIRMED ($R_\text{na}$ peaks at $\delta=5$) \\
$\beta_s$ & seeding scale & not swept & --- & Pending \\
\midrule
SS & quiescence gate & adversarial filter + commitment rate &
  P46+P47 & CONFIRMED (both effects) \\
\bottomrule
\end{tabular}
\end{table}

\section{Conclusion}

Experiment~A confirms that $\delta = W_\text{zone}/r_\text{wave}$ is the
operative third C$_2$ coordinate.  Zone encoding quality, as measured by the
non-adjacent/adjacent ratio, peaks at $\delta = 5$ and degrades toward random
at $\delta \approx 2$--$2.5$.  The practical minimum is $\delta \geq 5$ for
standard 4-zone grids with $r_\text{wave}=2$.

Experiment~B reveals a more complex interaction: at $T_\text{encode}$ near the
commitment epoch $t^*$, fidelity measures commitment rate rather than
adversarial resistance.  High-SS gates slow commitment (positive fidelity);
low-SS gates accelerate it (negative fidelity from sign flip).  This ties the
temporal gate invariant from Paper~46 directly to the commitment epoch dynamics
from Paper~26: \emph{the gate controls the timescale on which fieldM tracks
mid\_mem through the sign flip, and this timescale is the same as the
adversarial filtering timescale}.

The two-paper arc (Papers 46--47) closes the C$_2$ validation and opens a
natural next question: does there exist an optimal SS$^*$ as a function of
both adversarial pressure and proximity to $t^*$?  A joint sweep with
$T_\text{encode}$ varied would answer this directly.

\end{document}
